\documentclass[hyperref={unicode=true}]{beamer}
\usepackage[T2A]{fontenc}
\usepackage[utf8]{inputenx}
\usepackage[russian]{babel}

\usepackage{multicol}

\usepackage[pgf]{dot2texi}
\usepackage{tikz}
\usetikzlibrary{shapes, arrows}

\usepackage{listings}
\usepackage{graphicx}

\usepackage{comment}

\title{Лекция 5. Динамическая верификация}
\author{}
\date{}

\usetheme{Warsaw}

\AtBeginSection[] {
	\begin{frame}{Содержание}
		\tableofcontents[currentsection]
	\end{frame}
}
%\overfullrule=5pt

\begin{document}
	\begin{frame}{}
		\titlepage
	\end{frame}

    \begin{frame}{Цель лекции}
    Познакомиться с понятием динамической верификации
    и подробнее изучить способы написания мониторов ядра Linux.
    \end{frame}

    \section{Введение}

	\begin{frame}{Динамический анализ}
    \begin{itemize}
        \item valgrind, sanitizers -- поиск ошибок работы
        с памятью, ошибок в многопоточности, неопределенного
        поведения

        \item AFL (fuzzers) -- поиск входных данных, на которых
        происходят "крэши"

        \item очень удобные инструменты, но с фиксированными
        свойствами, которые они проверяют
	\end{itemize}
    \end{frame}

    \begin{frame}{Трудности тестирования}
    \begin{itemize}
        \item трудоемкость написания тестов с хорошим
        покрытием для крупных программ

        \item сложность воспроизведения ошибок
        в многопоточности
    \end{itemize}
    \end{frame}

    \begin{frame}{Идея динамической верификации}
    \begin{itemize}
       \item делать проверку прямо по ходу выполнения
       программы
       \item от тестирования берется возможность задать
       разные проверяемые свойства
       \item от динамического анализа берется
       мониторинг реального выполнения программы
    \end{itemize}
    \end{frame}

    \begin{frame}{Мониторинг вокруг нас}
    \begin{itemize}
        \item антивирусы
        \item защитники
        \item сетевые фильтры и анализаторы
    \end{itemize}
    \end{frame}

    \section{Терминология}

    \begin{frame}{Терминология}
    \begin{itemize}
        \item \emph{событие} (например, вход или выход
                из функции, достижение определенного адреса
                в программе и т.п.)
        \item \emph{трасса} --- последовательность событий
        \item выполняющаяся программа порождает трассу
        \item \emph{проверка трассы} на соответствие модели
        требований
    \end{itemize}
    \end{frame}

    \begin{frame}{Примеры моделей требований}
    \begin{itemize}
        \item регулярные выражения над алфавитом событий
        \item расширенные конечные автоматы
        \item темпоральные свойства
        \item машины Event-B
    \end{itemize}
    \end{frame}

    \begin{frame}{Пример свойства}
    \begin{itemize}
    \item \emph{Свойство}: перед вызовом \texttt{it.next()}
    должен быть вызов \texttt{it.hasNext()}, возвративший
    \texttt{true}, между которыми нет вызова \texttt{it.next()}
    \item \emph{События}: \texttt{next}, \texttt{hasNext(true)}
    \item \emph{Регулярное выражение}: \texttt{(hasNext(true) next)*}
    \end{itemize}
    \end{frame}

    \begin{frame}{Полный пример}
    \begin{itemize}
    \item Полный пример на Python, в котором проверяется, что
    программа успешно открывает только файлы в домашней директории.
    \item Находится по ссылке: \url{https://fmcourse.github.io/rv/open_home.py}
    \item Мониторинг делается утилитой \texttt{strace}
    \item Ее вывод перехватывается, анализируется и проверяется
    целевое свойство
    \end{itemize}
    \end{frame}

    \begin{frame}{Архитектура}
    \begin{itemize}
    \item Программа
    \item Монитор (наблюдает за ее выполнением, срабатывает во время
                событий)
    \item Модель требований (свойства трассы)
    \item Аниматор (проверяет трассу, например, воспроизводя ее на модели)
    \end{itemize}
    \end{frame}

    \begin{frame}{Вариации}
    \begin{itemize}
    \item офлайн-проверка / онлайн-проверка
    \item аниматор внутри монитора / аниматор отдельно
    \item инструментирование программы / динамическое подключение монитором
    \item тесты (заранее подготовленные сценарии) / реальная нагрузка
    \end{itemize}
    \end{frame}

    \section{eBPF}

    \begin{frame}{Что это}
    \begin{itemize}
    \item eBPF --- это инструментарий для разработки мониторов ядра Linux
    \item монитор ядра Linux --- это модуль ядра, который подключается
    к заранее заданным точкам в коде ядра Linux и выполняет заранее
    заданный код при достижении этих точек (\emph{eBPF-программа})
    \item чаще всего подключение происходит к системным вызовам
    \end{itemize}
    \end{frame}

    \begin{frame}{Особенности eBPF}
    \begin{itemize}
    \item можно перехватывать вход и выход в системный вызов
    (tracepoints)
    \item можно перехватывать вход и выход из произвольной экспортируемой
    функции ядра (kprobes)
    \item корректность eBPF-программы проверяется статически,
    а не во время выполнения (не требуется <<песочница>> во время
    ее выполнения в ядре)
    \item можно передавать данные в userspace через кольцевой буфер
    \end{itemize}
    \end{frame}

    \begin{frame}{Пример}
    \begin{itemize}
    \item Пример мониторинга ядра при помощи eBPF:
    \url{https://fmcourse.github.io/rv/ebpf.tgz}
    \item \texttt{syscall\_monitor.bpf.c} --- файл с eBPF-программами
    \item \texttt{monitor.c} --- утилита userspace
    \end{itemize}
    \end{frame}

    \begin{frame}{Пример (2)}
    \begin{itemize}
    \item Перехватываются системные вызовы
    \texttt{open}, \texttt{openat}, \texttt{exit}
    и \texttt{exit\_group}.
    \item В кольцевой буфер записывается структура
    \texttt{struct syscall\_event} с информацией о
    системном вызове
    \item В ней: номер системного вызова,
    pid процесса, запросившего системный вызов,
    и возвращаемое значение
    \end{itemize}
    \end{frame}

    \begin{frame}{Пример (3)}
    \begin{itemize}
    \item Утилита userspace запускает мониторинг
    программы, переданной в ее аргументах командной строки
    \item Во время работы программы она обрабатывает каждую
    структуру из кольцевого буфера
    \item Если эта структура про эту программу (совпадает pid),
    то печатает содержимое структуры
    \item Завершается по достижении системных вызовов
    \texttt{exit} и \texttt{exit\_group}
    \end{itemize}
    \end{frame}

    \begin{frame}{Жизненный цикл eBPF-программ}
    \begin{enumerate}
    \item Разработка kernelspace части: написание исходника (\texttt{.bpf.c})
    \item Компиляция: \texttt{clang -target bpf}
    \item Генерация скелетона: \texttt{bpftool gen skeleton file.bpf.o}
    \item Разработка userspace части: написание исходника и компиляция
    \item Загрузка eBPF-программ в память
    \item Присоединение их к точкам перехвата событий
    \item Исполнение
    \item Отключение и выгрузка из памяти
    \end{enumerate}
    \end{frame}

    \begin{frame}{Что не может делать eBPF-программа}
    \begin{itemize}
    \item использовать произвольно глобальные переменные
    \item вызывать функции (необходимо инлайнить все вызываемые
            функции)
    \item использовать операторы цикла, которые могут не остановиться
    \item если не сказано специально, то не могут вызывать
    действия, требующие переход процесса в состояние сна
    (например, обратиться к произвольному адресу с расчетом на то,
     в его отсутствие его страница будет подкачена в память)
    \item запрашивать, но не освобождать ресурсы
    \item использовать указатели на функции
    \end{itemize}
    \end{frame}

    \begin{frame}{Отладка}
    \begin{itemize}
    \item \texttt{bpf\_printk(...);} --- отладочная печать
    \item для просмотра нужно в реальном времени смотреть
    содержимое файла \texttt{/sys/kernel/debug/tracing/trace}
    (например, при помощи \texttt{cat})
    \end{itemize}
    \end{frame}

    \begin{frame}{Более сложные eBPF-программы}
    \begin{itemize}
    \item требуется запоминать аргументы системных вызовов
    и оказываемое воздействие (что поменялось), но аргументы
    недоступны на tracepoint sys\_exit $\Rightarrow$ нужно еще
    перехватывать tracepoint sys\_enter
    \item sys\_enter и sys\_exit --- это отдельные программы даже
    для одного и того же системного вызова, поэтому требуется
    глобальный map для передачи данных между ними
    \item оказываемое воздействие недоступно через аргументы
    системных вызовов (например, inode созданного файла), но
    их можно получить, перехватывая аргументы функций,
    вызываемых при выполнении системного вызова
    \end{itemize}
    \end{frame}

    \begin{frame}{Более сложные eBPF-программы (2)}
    \begin{itemize}
    \item если системных вызовов надо обрабатывать много, то
    получается много eBPF-программ, их загрузка может занимать
    значительное время
    \item eBPF-программа обрабатывает события от всех процессов,
    поэтому требуется фильтрация событий внутри eBPF-программ,
    иначе кольцевой буфер переполнится
    \item отладка всё ещё непростая
    \end{itemize}
    \end{frame}

    \begin{frame}{Вперед!}
    Изучите более сложные eBPF-программы.
    \end{frame}

\end{document}

